\begin{lemma}
For every $\eta>0$ there is $D_0$ such that, whenever $D\ge D_0$,
$\mathcal H$ is a $3$-uniform $3$-simple hypergraph with maximum degree at
most $D$, and $e\in E(\mathcal H)$, one has
\[
b_{L(\mathcal H),3D}(e)\le 1-\frac19+\frac1{3^5}+\eta .
\]
\end{lemma}
\begin{proof}
Fix $\eta>0$. Choose a uniform integer threshold large enough for
\noderef{ExtremalBParameterSaturationThreeSimple},
\noderef{ThreeUniformThreeSimplePairUpperCutoff}, and
\noderef{ThreeUniformThreeSimpleT2UniformEstimate} with error $\eta$,
and at least $100$.

Let $D$ exceed this threshold and let $(\mathcal H,e)$ be admissible.
First apply \noderef{FiniteHypergraphClassRelabeling} with $k=t=3$ and
$M=3D$. Its scale hypotheses hold since $D>0$. This gives finite vertex
and edge types, class membership, and exactly the same local parameter.
Apply \noderef{KUniformKSimpleLocalPatternExtremalReduction} at $k=3$
to this relabeled hypergraph. It supplies a representative
$(\mathcal F_m,f_m)$ with local parameter at least the original one
and at least that of every competitor in $H(3,3,D)$. Its vertex type
need not be finite. Apply \noderef{FiniteHypergraphClassRelabeling}
again to this representative. The resulting $(\mathcal F,f)$ has
finite vertex and edge types and the same parameter as
$(\mathcal F_m,f_m)$, hence retains both inequalities and global
maximality. In Lean both applications may land in fixed small universes;
the first also ensures that the extremal reduction and all subsequent
competitors use those universes.

Apply \noderef{ExtremalBParameterSaturationThreeSimple} to this
maximizer. Every vertex of $f$ has degree $D$ and $f$ has exactly
$3(D-1)$ neighbors. Together with $D\ge100$, finite types, class
membership, and maximality these are exactly the assumptions of
\noderef{ThreeUniformThreeSimpleSetupConstruction}, which supplies a
setup $S$ on the same $\mathcal F,f$. Thus no setup-dependent estimate
is used to produce the setup.

Write $P,T,b$ for its pair count, triple count, and parameter.
By \noderef{ThreeUniformThreeSimplePairLowerBound} and
\noderef{ThreeUniformThreeSimplePairUpperCutoff},
\[
 D^2-9D+Y\le P\le1.031D^2-9D .
\]
Define $\delta=(P+9D)/D^2-1$. Since $D>0$, this gives
$P=(1+\delta)D^2-9D$ exactly. Using $Y\ge0$ from
\noderef{ThreeUniformThreeSimpleLocalSetup} in the lower bound and
dividing both bounds by $D^2>0$ gives $0\le\delta\le0.031$.

The cutoff and pair equality now allow
\noderef{ThreeUniformThreeSimpleSmallXs} to give
$\mathrm{totalXs}\le(8/5)\delta D$.
By \noderef{ThreeUniformThreeSimpleT1IncidenceBound},
\[
 T_1\le |V_1|D^2\le\mathrm{totalXs}\,D^2
       \le(8/5)\delta D^3.
\]
The uniform estimate \noderef{ThreeUniformThreeSimpleT2UniformEstimate}
gives
\[
 T_2\le(1/9+1.33\delta+\eta)D^3 .
\]
That estimate includes the positive auxiliary epsilon slack needed when
$\delta=0$; its threshold is uniform in the setup and $\delta$.
Since $T=T_1+T_2$ by the setup, and $\delta\ge0$,
\[
 T\le(1/9+2.93\delta+\eta)D^3
   \le(1/9+3\delta+\eta)D^3.
\]

Finally the exact parameter identity in
\noderef{ThreeUniformThreeSimpleNeighborhoodComplementBounds} gives
\[
\begin{aligned}
 b&=\frac{3(D-1)}{3D}-\frac{P}{(3D)^2}+\frac{T}{(3D)^3}\\
 &\le\frac{3(D-1)}{3D}
 -\frac{(1+\delta)D^2-9D}{(3D)^2}
 +\frac{(1/9+3\delta+\eta)D^3}{(3D)^3}\\
 &=1-\frac19+\frac1{3^5}+\frac{\eta}{27}
 \le1-\frac19+\frac1{3^5}+\eta.
\end{aligned}
\]
The equality cancels $-1/D$ against $+1/D$ and
$-\delta/9$ against $+\delta/9$. The last inequality uses $\eta>0$.
The extremal comparison and both relabeling equalities transfer this
bound to the original $(\mathcal H,e)$, as required.
\end{proof}
